\DocumentMetadata{
  pdfversion=2.0,
  pdfstandard=ua-2,
  lang=en-US,
  tagging=on,
  tagging-setup={math/setup=mathml-SE,tabsorder=structure}
}
\documentclass[11pt,letterpaper]{article}
\usepackage[margin=0.68in,includefoot]{geometry}
\usepackage{newcomputermodern}
\usepackage{amsmath,amscd,mathtools}
\usepackage{tikz,adjustbox}
\providecommand{\square}{\mdlgwhtsquare}
\usepackage[hidelinks]{hyperref}
\hypersetup{
  pdftitle={Topology First-Year Exam, August 2015, Part 2},
  pdfauthor={University of Florida Department of Mathematics},
  pdfsubject={Graduate mathematics examination},
  pdfdisplaydoctitle=true,
  bookmarksopen=true
}
\setcounter{secnumdepth}{0}
\setlength{\parindent}{0pt}
\setlength{\parskip}{0.28em}
\setlength{\emergencystretch}{3em}
\setlength{\leftmargini}{2.15em}
\setlength{\leftmarginii}{2.65em}
\setlength{\leftmarginiii}{3em}
\renewcommand{\labelenumi}{\textbf{\theenumi.}}
\pagestyle{empty}
\raggedbottom
\allowdisplaybreaks
\newenvironment{examproblems}[1]{%
  \begin{enumerate}%
  \setcounter{enumi}{#1}%
  \setlength{\topsep}{0.35em}%
  \setlength{\partopsep}{0pt}%
  \setlength{\itemsep}{0.48em}%
  \setlength{\parsep}{0.18em}%
}{\end{enumerate}}
\newenvironment{examsubparts}{%
  \begin{enumerate}%
  \setlength{\topsep}{0.18em}%
  \setlength{\partopsep}{0pt}%
  \setlength{\itemsep}{0.16em}%
  \setlength{\parsep}{0.1em}%
}{\end{enumerate}}
\newenvironment{examinstructions}{%
  \par\begingroup\itshape
}{%
  \par\endgroup\vspace{0.18em}
}
\makeatletter
\renewcommand\section{\@startsection{section}{1}{0pt}%
  {-1.25ex plus -.25ex minus -.1ex}%
  {0.55ex plus .1ex}%
  {\normalfont\large\bfseries}}
\renewcommand{\@maketitle}{%
  \newpage\null\vskip -1em%
  {\footnotesize\bfseries
    \href{https://www.ufl.edu/}{University of Florida}, \href{https://math.ufl.edu/}{Department of Mathematics}\par}%
  \vskip -0.5em%
  \tagstructbegin{tag=Title}%
    {\LARGE\bfseries\href{https://gma.math.ufl.edu/exams/topology-first-year/}{Topology First-Year Exam}, August 2015, Part 2\par}%
  \tagstructend%
  \vskip 0.25em%
  \tagmcbegin{artifact}%
    \hrule height 1pt%
  \tagmcend%
  \vskip 1em%
}
\makeatother
\title{Topology First-Year Exam, August 2015, Part 2}
\author{}
\date{}
\begin{document}
\maketitle
\thispagestyle{empty}
\begin{examinstructions}
Answer all questions and work all problems.
\end{examinstructions}
\begin{examproblems}{0}
\item
Let \(n>1\). Suppose that \(f,g:X\to S^n\) are continuous such that \(f(x)\) and \(g(x)\) are not antipodal for all \(x\in X\). Show that~\(f\) and~\(g\) are homotopic.
\item
Show that \(\pi_1(S^n,1)=0\) for all \(n>1\).
\item
Show that \(\pi_1(S^1,1)=\mathbb{Z}\).
\item
Show that the disk, \(D^2=\{z\in\mathbb{C}\mid \lVert z\rVert\leq 1\}\) has the fixed point property.
\item
Consider the matrix \(M=\begin{bmatrix}2&3\\1&1\end{bmatrix}\). This represents a group homomorphism \(M:\mathbb{Z}^2\to\mathbb{Z}^2\). Show that there is a map \(f:\mathbb{T}^2\to\mathbb{T}^2\) such that \(f_*=M\), where \(f_*:\pi_1(\mathbb{T}^2)\to\pi_1(\mathbb{T}^2)\) is the homomorphism induced by~\(f\) on the fundamental group.
\item
Consider three loops \(f,g,h:[0,1]\to(X,x_0)\). Show that \((f\circ g)\circ h\) is homotopic to \(f\circ(g\circ h)\), where~\(\circ\) denotes concatenation of loops.
\item
Let \(n>1\). Let \(P^n=S^n/D\), where~\(D\) identifies~\(x\) with its antipode for all \(x\in S^n\). Show that \(\pi_1(P^n,x_0)\cong\mathbb{Z}_2\).
\item
Describe a CW-complex~\(X\) such that \(\pi_1(X,x_0)\cong\mathbb{Z}_3\).
\item
Suppose that \(\pi_1(X,x_0)\cong\mathbb{Z}_2\) and that \(\pi_1(Y,y_0)\cong\mathbb{Z}_5\). Show that \(\pi_1(X\times Y,(x_0,y_0))\cong\mathbb{Z}_{10}\).
\item
Suppose that~\(X\) and~\(Y\) are connected. Show that \(X\times Y\) is connected.
\item
Let \(n>1\). Let \(A\subset\mathbb{R}^n\) be a countable set. Show that \(\mathbb{R}^n\setminus A\) is arcwise connected.
\item
Let \(X=S^1\vee S^1\). Show that \(\pi_1(X)\cong\mathbb{Z}*\mathbb{Z}\).
\item
State the following theorems.

\par
Seifert–van Kampen Theorem.

\par
The Urysohn Metrization Theorem

\par
The Urysohn Lemma

\par
The Tietze Extension Theorem

\par
The Brouwer Fixed Point Theorem

\par
The Fundamental Theorem of Algebra

\par
The Jordan Curve Theorem

\par
The Arcwise Connectedness Theorem

\par
The Hahn–Mazurkiewicz Theorem
\end{examproblems}
\end{document}
